% ==========================================================================
% OG-Core x CLEWS/OSeMOSYS — THE COUPLING MECHANISMS: A REFERENCE.
% One entry per channel: the intuition, then the technical form, then the
% variables. Companion to the deck (mechanisms.tex) and the paper (../paper);
% numbers are the committed M=8 Philippine golden record.
% Build: latexmk -pdf mechanisms-reference.tex
% ==========================================================================
\documentclass[11pt]{article}
\usepackage[T1]{fontenc}
\usepackage[margin=1in]{geometry}
\usepackage{amsmath,amssymb}
\usepackage{microtype}
\usepackage{enumitem}
\usepackage{parskip}
\usepackage{xcolor}
% OG-Core x CLEWS — shared color palette.
% Single source of truth, mirrored from ogclews_link/style.py (the matplotlib figure
% theme) so the TikZ diagrams, the Beamer deck, and the data figures all share one
% editorial, colorblind-safe palette. \input this from any preamble; it only defines
% colors (no package loads), so it is safe in both standalone and beamer documents.

\definecolor{ink}{HTML}{222222}      % primary text / heavy rules
\definecolor{sub}{HTML}{555555}      % secondary text
\definecolor{mute}{HTML}{888888}     % tertiary / source lines
\definecolor{grid}{HTML}{E6E6E6}     % faint rules, fills

% the two models + policy — the load-bearing semantic colors
\definecolor{ogc}{HTML}{0F5499}      % OG-Core      (FT oxford blue)
\definecolor{clews}{HTML}{0D7680}    % CLEWS/OSeMOSYS (teal)
\definecolor{policy}{HTML}{E3120B}   % policy lever (Economist red)

% accents / signed values (match style.py)
\definecolor{claret}{HTML}{990F3D}   % kicker rule
\definecolor{gain}{HTML}{2166AC}     % positive change (blue)
\definecolor{loss}{HTML}{B2182B}     % negative change (red)
\definecolor{amber}{HTML}{E69F00}    % categorical 4
\definecolor{purple}{HTML}{7A6FAC}   % categorical 5
\definecolor{slate}{HTML}{7F8C8D}    % categorical 6 / neutral

% sequential blue ramp (ordered income groups, poorest -> richest)
\definecolor{seqA}{HTML}{C6DBEF}
\definecolor{seqB}{HTML}{9ECAE1}
\definecolor{seqC}{HTML}{6BAED6}
\definecolor{seqD}{HTML}{4292C6}
\definecolor{seqE}{HTML}{2171B5}
\definecolor{seqF}{HTML}{08519C}
\definecolor{seqG}{HTML}{08306B}

\usepackage[colorlinks=true,linkcolor=ogc,citecolor=ogc,urlcolor=clews]{hyperref}

% ---- notation (consistent with the deck and paper) ---------------------------
\newcommand{\OGCore}{\textsc{OG-Core}}
\newcommand{\CLEWS}{\textsc{CLEWS}}
\newcommand{\OSeMOSYS}{\textsc{OSeMOSYS}}
\newcommand{\tauc}{\tau^{c}}
\newcommand{\taub}{\tau^{b}}
\newcommand{\alphaI}{\alpha_{I}}
\newcommand{\alphaT}{\alpha_{T}}
\newcommand{\Kg}{K_{g}}
\newcommand{\gammam}{\gamma_{m}}
\newcommand{\gammag}{\gamma_{g}}
\newcommand{\rp}{r_{p}}

% ---- per-mechanism template --------------------------------------------------
\newcommand{\dir}[1]{\nolinebreak\hfill{\normalfont\small\color{mute}[#1]}}
\newcommand{\Intuition}{\par\smallskip\noindent{\color{ogc}\bfseries Intuition.}\ }
\newcommand{\Technical}{\par\smallskip\noindent{\color{clews}\bfseries Technical.}\ }
\newcommand{\Variables}{\par\smallskip\noindent{\color{sub}\bfseries Variables.}\ }

\setlist[itemize]{leftmargin=1.4em, itemsep=1pt, topsep=2pt}

\title{\vspace{-1.2cm}\bfseries OG-Core\,$\times$\,CLEWS --- The Coupling Mechanisms}
\author{}
\date{A reference: the intuition, then the technical form, for each channel.\\
\small Numbers are the committed M=8 Philippine golden record.}

\begin{document}
\maketitle
\vspace{-0.8cm}

% ==========================================================================
\section{The coupling in brief}

\OGCore{} and \CLEWS{} are complementary. \CLEWS{}/\OSeMOSYS{} is a least-cost energy
optimiser --- rich in technology, dispatch, and emissions, but blind to behaviour: its
demand is a fixed input, and it has no households, government budget, or distribution.
\OGCore{} is an overlapping-generations general-equilibrium economy --- rich in
households (by age $s$ and lifetime income $j$), firms ($M$ industries), and a
government that closes its budget --- but blind to energy: energy is not a priced input
to firm production. They are joined by a \emph{staged soft link} that exchanges a few
signals; a single hard model is ruled out because one steps through time period by
period while the other plans all years at once.

Three objects organise every mechanism:
\begin{itemize}
  \item \textbf{Signal} --- a quantity read from a \emph{solved} \CLEWS{} scenario: a
    fuel's marginal cost, the capital a build-out needs, emissions, an implied carbon price.
  \item \textbf{Wedge} --- the \OGCore{} parameter that carries it: the consumption tax
    $\tauc$, an industry productivity $Z_m$, public investment $\alphaI$, mortality
    $\rho(s,t)$, effective labour $e(t,s,j)$.
  \item \textbf{Channel} --- one signal mapped to one wedge, with a guard that fires
    when the mapping is large or ill-posed.
\end{itemize}

The load-bearing signal is the \textbf{energy price}: \CLEWS{}'s marginal cost of
supplying one more unit. Fed to \OGCore{}, it
does two things \CLEWS{} cannot --- it lets demand respond to price, and it reveals who
bears the cost. Because \OGCore{} firms have no energy input, the energy-price channel represents an
electricity-price change through electricity's \emph{two} uses, combined in one
transmission: a cost-push on the industries that buy electricity as an input (the
contractionary, wage-bearing side), and a recycled consumption wedge on households (the
regressive cost-of-living side). The input side is a reduced-form proxy; the
fully-structural version --- energy as a priced production input --- is not yet built.

Throughout, a hat denotes a reform/baseline ratio ($\hat x \equiv x^{r}/x^{b}$), and
$i_e$ / $m_e$ are the energy \emph{good} / \emph{industry} (Philippines: $m_e=2$).

% ==========================================================================
\section{The mechanisms}

The channels split by direction: four carry \CLEWS{} signals into \OGCore{}, two are
policy levers applied once to both models, and two send \OGCore{} results back to
\CLEWS{}.

\subsection{Energy price \dir{CLEWS $\to$ OG}}
\Intuition An electricity-price change hits the economy in two places, because
electricity is used in two ways. Industries buy most of it as an \emph{input} (in the
Philippine accounts about $73\%$), so a higher price raises their costs --- and, through
them, wages and output. Households buy the rest \emph{directly} (about $25\%$), so it
raises the cost of living, hardest for the poor (energy is a larger share of a poor
budget). The channel is one complete transmission that carries both.
\Technical The reform/baseline electricity-price ratio $g$ is applied through both uses.
\emph{(i) Firm input --- a cost-push:} each industry $j$'s productivity is cut by
electricity's share $\phi_j$ of its input costs (from the SAM),
$Z_j \mapsto Z_j/(1+\phi_j(g-1))$, so its output price rises by $\approx\phi_j(g-1)$ ---
the contractionary, wage-bearing side. \emph{(ii) Household consumption --- a wedge:} a
recycled consumption-tax wedge on the energy good, scaled to electricity's value-share
$\sigma$ of that good,
\[
  1+\tau^{c,r}_{i_e} \;=\; \bigl(1+\sigma(g-1)\bigr)\bigl(1+\tau^{c,b}_{i_e}\bigr)
\]
--- the regressive cost-of-living side. Electricity's own self-use is zeroed in the
cost-push so it is not double-counted by the wedge. The input side is a reduced-form
proxy ($\OGCore{}$ has no energy in production; the structural version is a real
inter-industry use matrix, unbuilt). The price $g$ is \CLEWS{}'s ``auto'' source --- a
cost-of-electricity index if available, else the \OSeMOSYS{} commodity-balance marginal cost; for
the Philippines it is near-flat, so the whole channel is small.
\Variables $g=\pi^{r}/\pi^{b}$, the reform/baseline electricity-price ratio; $\phi_j$,
electricity's share of industry $j$'s input costs; $\sigma$, electricity's value-share of
the household energy good $i_e$; $Z_j$, industry $j$'s productivity; $\tau^{c}_{i_e}$, the
consumption-tax wedge on the energy good ($r,b$ = reform, baseline).

\subsection{Carbon price \dir{policy --- both models}}
\Intuition One carbon price, chosen by policy, applied once on both sides: in the energy
model it tilts the build toward cleaner technology; in the economy it raises the price of
(household) energy like a tax. It is counted once --- never also inferred back from the
energy model's own implied carbon cost.
\Technical A price $cp$ (US\$/tCO$_2$) becomes an add-on to the energy-good wedge in
\OGCore{}, and an emissions-penalty path in \CLEWS{}:
\[
  \Delta\tauc_{i_e} \;=\; \frac{cp \,\cdot\, \text{deflator}\,\cdot\,\text{carbon intensity}}{p^{b}_{i_e}}.
\]
Because \OGCore{} firms buy no energy, only \emph{household} energy carbon is priced on
the macro side; industrial carbon is priced in \CLEWS{}. Revenue is recyclable as
transfers ($\alphaT$). Standalone $Y_{ss}=-0.040\%$; the \emph{level} rides on the
uncalibrated unit/deflator bridge, so treat the magnitude as illustrative.
\Variables $cp$, the carbon price; the deflator maps \CLEWS{} money to the \OGCore{}
numéraire (currently $1$, uncalibrated); carbon intensity, tCO$_2$ per unit of the energy
good; $\alphaT$, the transfer share used to recycle revenue.

\subsection{Public investment \dir{CLEWS $\to$ OG}}
\Intuition Building the grid (transmission and distribution) is a public cost. If
debt-financed it competes with everyone else for savings, pushing up the interest rate
and crowding out private investment; but public capital is also productive, so it lifts
output. The net sign depends on which dominates.
\Technical The public grid-capex increment (as a share of GDP) is added to the
public-investment share $\alphaI$ on a finite ramp; it accumulates into public capital,
which is productive in the firm's technology:
\[
  K_{g,t+1} = (1-\delta)\,K_{g,t} + I_{g,t}, \qquad I_{g,t} = \alpha_{I,t}\,Y_t .
\]
This is where crowding-out legitimately lives. In the Philippine scenario the extra grid
spend is $\approx 0$ (the transition is generation-side), so the channel is inert here
($Y_{ss}=+0.000\%$); its mechanics scale with any non-zero spend.
\Variables $\alphaI$, the public-investment share of output; $\Kg$, the public capital
stock; $I_g$, public investment; $\delta$, depreciation; $\gammag$, public capital's
share in production.

\subsection{Capital intensity \dir{CLEWS $\to$ OG}}
\Intuition A clean power fleet is capital-heavy (much up-front capex, little fuel). The
tempting reading --- ``so it pulls capital in'' --- is wrong here. Raising the energy
sector's capital \emph{share} with productivity fixed is like making the sector cheaper
to run: competition drives the electricity price down, and because electricity demand
barely responds, the sector's revenue shrinks and it \emph{sheds} capital. So this is a
factor-income / energy-price lever, not a capital draw-in.
\Technical The reform/baseline ratio of the power fleet's capital-cost share scales the
energy industry's capital exponent $\gammam$ (Philippines: $+12.2\%$). Under
Cobb--Douglas technology ($\varepsilon=1$) capital's revenue share \emph{is} $\gammam$,
so firm capital demand is $K_m=\gammam\,p_m Y_m/\rho_m$. The lever moves only $\gammam$,
so in the reform/baseline ratio the cost of capital $\rho_m$ and the return $r$ cancel,
leaving an exact identity:
\[
  \hat K_E \;=\; \hat\gamma_E \cdot \hat p_E \cdot \hat Y_E .
\]
Measured (energy industry): $\hat\gamma_E=+12.2\%$, $\hat p_E=-42.8\%$,
$\hat Y_E=+20.9\%$, so $\hat K_E=-22.4\%$ (the identity reproduces it to $0.03\%$).
\Variables $\gammam$, industry $m$'s capital exponent (its share of value added);
$\rho_m=(r+\delta-\taub_m\delta^{\tau}_m-\tau^{\mathrm{inv}}_m\delta)/(1-\taub_m)$, the
cost of capital; $p_m,Y_m$, the industry's price and output; a hat is a reform/baseline
ratio; $E$ marks the energy industry.

\subsection{Investment incentive \dir{policy}}
\Intuition To actually pull capital \emph{into} energy, lower its cost of capital --- for
instance with an investment tax credit. That makes energy investment cheaper at the going
interest rate, so capital reallocates in. This is a genuinely different lever from capital
intensity: capital intensity changes the factor \emph{mix} (the exponent); the credit
changes the \emph{cost of capital}. They move energy capital in opposite directions, so
they are not two views of the same thing.
\Technical An investment tax credit $\tau^{\mathrm{inv}}_m$ on the energy industry enters
the cost of capital directly, where the capital exponent is \emph{absent}:
\[
  \rho_m=\frac{r+\delta-\taub_m\delta^{\tau}_m-\tau^{\mathrm{inv}}_m\delta}{1-\taub_m}.
\]
A higher $\tau^{\mathrm{inv}}_m$ lowers $\rho_m$, so capital demand $K_m=\gammam p_m Y_m/\rho_m$
shifts out at the prevailing $r$. It is a defined lever, not part of the Philippine
reference run.
\Variables $\tau^{\mathrm{inv}}_m$, the investment tax credit rate; $\taub_m$, the
business-income tax; $\delta$, economic depreciation; $\delta^{\tau}_m$, tax depreciation;
$\rho_m$, the cost of capital.

\subsection{Health \dir{CLEWS $\to$ OG}}
\Intuition Cleaner air saves lives and --- more important for output --- makes
working-age people healthier and more productive. But a $7.7\%$ cut in power-sector
PM$_{2.5}$ is not $7.7\%$ fewer deaths: energy is one pollution source among many, and
deaths fall less than one-for-one as the air clears, so only $\sim$8\% of the cut reaches
mortality. The deaths avoided skew elderly (retirees), so the long-run output effect is
$\approx 0$; the real macro gain is the working-age productivity channel, which shows up
along the transition.
\Technical The reform/baseline PM$_{2.5}$ emissions change $\Delta\text{em}$ scales the
country's Global Burden of Disease attributable deaths by a calibrated dose-response $M$:
\[
  \Delta\text{deaths} \;=\; D_{\mathrm{PM}} \times M \times \Delta\text{em}
  \;=\; 43{,}951 \times 0.082 \times (-0.077) \;\approx\; -279 .
\]
\emph{Mortality} shocks the age-specific rate $\rho(s,t)$ to hit that signed target and
recomputes the population; \emph{morbidity} scales the effective-labour profile
$e(t,s,j)$ by a working-age disability shape. Result: $Y_{ss}=-0.003\%$ but $+0.19\%$
near-term. The $\times M$ is what separates $-279$ from the naïve $-3{,}400$.
\Variables $\Delta\text{em}$, the reform/baseline PM$_{2.5}$ emissions change;
$D_{\mathrm{PM}}$, GBD ambient-PM$_{2.5}$ attributable deaths ($43{,}951$, Philippines);
$M\approx0.082$, the dose-response (energy mass-share $\times$ concentration-response
elasticity); $\rho(s,t)$, mortality by age $s$ and time $t$; $e(t,s,j)$, effective labour
by time, age, and income group $j$.

\subsection{Discount rate \dir{OG $\to$ CLEWS}}
\Intuition The economy has an equilibrium return on capital --- its own cost of capital.
Handing that to the energy planner means it discounts future costs at a rate the economy
actually implies, rather than an arbitrary assumption. A lower return favours
capital-heavy, long-lived assets (renewables).
\Technical After the reform solves, \OGCore{}'s equilibrium market return $\rp$ (a
multi-period mean) is written to \CLEWS{}'s planning discount rate. It is emitted
\emph{after} the solve and is one-way: it changes the energy build, whose prices feed back
to \OGCore{} only once the loop iterates. Inert as a standalone reform.
\Variables $\rp$, the market portfolio return (\OGCore{}'s equilibrium cost of capital);
the \CLEWS{} discount rate it sets.

\subsection{Energy demand \dir{OG $\to$ CLEWS}}
\Intuition A larger or richer economy needs more energy. Once the economy solves, its
activity scales the demand the energy system must meet. In a single pass this barely
changes, so the channel is latent --- it only bites once the two models iterate.
\Technical The reform/baseline activity ratio (industry output $Y_m$ or consumption
$C_i$, raised to an elasticity) scales \CLEWS{}'s specified annual demand:
$\text{demand}^{r} = (\hat Y_m)^{\varepsilon_D}\,\text{demand}^{b}$. Emitted after the
solve; one-way. In one pass $\hat Y_m\approx 1$, so it is inert by construction.
\Variables $Y_m$, industry output; $C_i$, consumption of good $i$; $\varepsilon_D$, the
demand-scaling elasticity; the \CLEWS{} specified annual demand it scales.

\end{document}
